\documentclass{amsart}
% For dealing with references we use the comment environment
\usepackage{verbatim}
\newenvironment{reference}{\comment}{\endcomment}
%\newenvironment{reference}{}{}
\newenvironment{slogan}{\comment}{\endcomment}
\newenvironment{history}{\comment}{\endcomment}

% For commutative diagrams we use Xy-pic
\usepackage[all]{xy}

% We use 2cell for 2-commutative diagrams.
\xyoption{2cell}
\UseAllTwocells

% We use multicol for the list of chapters between chapters
\usepackage{multicol}

% This is generally recommended for better output
\usepackage{lmodern}
\usepackage[T1]{fontenc}

% For cross-file-references

% Package for hypertext links:
\usepackage{hyperref}

% For any local file, say "hello.tex" you want to link to please

% Theorem environments.
%
\theoremstyle{plain}
\newtheorem{theorem}[subsection]{Theorem}
\newtheorem{proposition}[subsection]{Proposition}
\newtheorem{lemma}[subsection]{Lemma}

\theoremstyle{definition}
\newtheorem{definition}[subsection]{Definition}
\newtheorem{example}[subsection]{Example}
\newtheorem{exercise}[subsection]{Exercise}
\newtheorem{situation}[subsection]{Situation}

\theoremstyle{remark}
\newtheorem{remark}[subsection]{Remark}
\newtheorem{remarks}[subsection]{Remarks}

\numberwithin{equation}{subsection}

% Macros
%
\def\lim{\mathop{\mathrm{lim}}\nolimits}
\def\colim{\mathop{\mathrm{colim}}\nolimits}
\def\Spec{\mathop{\mathrm{Spec}}}
\def\Hom{\mathop{\mathrm{Hom}}\nolimits}
\def\Ext{\mathop{\mathrm{Ext}}\nolimits}
\def\SheafHom{\mathop{\mathcal{H}\!\mathit{om}}\nolimits}
\def\SheafExt{\mathop{\mathcal{E}\!\mathit{xt}}\nolimits}
\def\Sch{\mathit{Sch}}
\def\Mor{\mathop{\mathrm{Mor}}\nolimits}
\def\Ob{\mathop{\mathrm{Ob}}\nolimits}
\def\Sh{\mathop{\mathit{Sh}}\nolimits}
\def\NL{\mathop{N\!L}\nolimits}
\def\CH{\mathop{\mathrm{CH}}\nolimits}
\def\proetale{{pro\text{-}\acute{e}tale}}
\def\etale{{\acute{e}tale}}
\def\QCoh{\mathit{QCoh}}
\def\Ker{\mathop{\mathrm{Ker}}}
\def\Im{\mathop{\mathrm{Im}}}
\def\Coker{\mathop{\mathrm{Coker}}}
\def\Coim{\mathop{\mathrm{Coim}}}

% Boxtimes
%
\DeclareMathSymbol{\boxtimes}{\mathbin}{AMSa}{"02}

%
% Macros for moduli stacks/spaces
%
\def\QCohstack{\mathcal{QC}\!\mathit{oh}}
\def\Cohstack{\mathcal{C}\!\mathit{oh}}
\def\Spacesstack{\mathcal{S}\!\mathit{paces}}
\def\Quotfunctor{\mathrm{Quot}}
\def\Hilbfunctor{\mathrm{Hilb}}
\def\Curvesstack{\mathcal{C}\!\mathit{urves}}
\def\Polarizedstack{\mathcal{P}\!\mathit{olarized}}
\def\Complexesstack{\mathcal{C}\!\mathit{omplexes}}
% \Pic is the operator that assigns to X its picard group, usage \Pic(X)
% \Picardstack_{X/B} denotes the Picard stack of X over B
% \Picardfunctor_{X/B} denotes the Picard functor of X over B
\def\Pic{\mathop{\mathrm{Pic}}\nolimits}
\def\Picardstack{\mathcal{P}\!\mathit{ic}}
\def\Picardfunctor{\mathrm{Pic}}
\def\Deformationcategory{\mathcal{D}\!\mathit{ef}}

\setlength{\emergencystretch}{3em}
\begin{document}
\title[Stacks Project, Tag 0H41]{486 --- Finite-presentation scheme-theoretic images commute with arbitrary base change on a finite closed stratification after a finite-type thickening of the base}
\maketitle
\noindent Copyright (C) 2005--2025 Johan de Jong. Stacks Project authors. Source: \href{https://stacks.math.columbia.edu/tag/0H41}{Tag 0H41}.

\noindent Editorial difficulty note (not author text). Difficulty H2: Noetherian approximation and induction must produce a finite stratification with finite-type defining ideals over the original quasi-compact quasi-separated base. On every generic stratum, cokernel flatness and simultaneous flatness of the two algebras ensure kernel formation commutes with all tensor products, then affine-local image descriptions globalize this condition..

\noindent Editorial provenance note (not author text). History: excerpt from revision \texttt{a0444\allowbreak 6e57e\allowbreak c1fbc\allowbreak 252a8\allowbreak 71afc\allowbreak ec775\allowbreak 2fb28\allowbreak 07b14}, more-morphisms.tex, lines 15775--15866. Reference presentation was adapted; original-author context and core support blocks were retained or added. The mathematical wording is unchanged. Licensed under GNU FDL 1.2 or later; no invariant sections or cover texts. See ../licenses/COPYING. Context blocks reproduce author definitions and supporting results; they are not additional counted problems.

\noindent
Let $f : X \to Y$ be a morphism of schemes over a base scheme $S$.
Let $Z \subset Y$ be the scheme theoretic image of $f$, see
Morphisms, Section \href{https://stacks.math.columbia.edu/tag/01R5}{section scheme theoretic image (Tag 01R5)}.
Let $g : S' \to S$ be a morphism of schemes and let
$f' : X \times_S S' \to Y \times_S S'$ be the base change of $f$ by $g$.
It is not always true that $Z \times_S S' \subset Y \times_S S'$
is the scheme theoretic image of $f'$.
Let us say that {\it formation of the scheme theoretic image of $f/S$
commutes with arbitrary base change} if for every $g$ as above
the scheme theoretic image of $f'$ is equal to $Z \times_S S'$.


\begin{lemma}
\label{lemma-base-change-scheme-theoretic-image}
Let $S$ be a quasi-compact and quasi-separated scheme.
Let $f : X \to Y$ be a morphism of schemes over $S$ with
both $X$ and $Y$ of finite presentation over $S$.
Then there exists a $t \geq 0$ and closed subschemes
$$
S \supset S_0 \supset S_1 \supset \ldots \supset S_t = \emptyset
$$
with the following properties:
\begin{enumerate}
\item $S_i \to S$ is defined by a finite type ideal sheaf,
\item $S_0 \subset S$ is a thickening, and
\item with $T_i = S_i \setminus S_{i + 1}$ and $f_i$ the base
change of $f$ to $T_i$ we have:
formation of the scheme theoretic image of $f_i/T_i$
commutes with arbitrary base change (see discussion above the lemma).
\end{enumerate}
\end{lemma}

\begin{proof}
We can find a commutative diagram
$$
\xymatrix{
X \ar[d] \ar[r] & Y \ar[d] \ar[r] & S \ar[d] \\
U \ar[r] & V \ar[r] & W
}
$$
with cartesian squares such that $U$, $V$, $W$
are of finite type over $\mathbf{Z}$. Namely, first write $S$
as a cofiltered limit of finite type schemes over $\mathbf{Z}$
with affine transition morphisms
using Limits, Proposition \href{https://stacks.math.columbia.edu/tag/01ZA}{proposition approximate (Tag 01ZA)}
and then descend the morphism $X \to Y$ using
Limits, Lemma \href{https://stacks.math.columbia.edu/tag/01ZM}{lemma descend finite presentation (Tag 01ZM)}.
This reduces us to the case discussed in the next paragraph.

\medskip\noindent
Assume $S$ is Noetherian.  In this case every quasi-coherent ideal
is of finite type, hence we do not have to check the condition that
$S_i$ is cut out by a finite type ideal. Set $S_0 = S_{red}$ equal
to the reduction of $S$. Let $\eta \in S_0$ be a generic point
of an irreducible component of $S_0$. By Noetherian induction on
the underlying topological space of $S_0$, we may assume the result
holds for any closed subscheme of $S_0$ not containing $\eta$.
Thus it suffices to show that there exists an open neighbourhood
$U_0 \subset S_0$ such that the base change $f_0$ of $f$ to $U_0$
has property (3).

\medskip\noindent
Let $R$ be a Noetherian domain. Let $f : X \to Y$ be a morphism
of finite type schemes over $R$. By the discussion in the previous paragraph
it suffices to show that after replacing $R$ by $R_g$ for some $g \in R$
nonzero and $X$, $Y$ by their base changes to $R_g$,
formation of the scheme theoretic image of $f/R$ commutes
with arbitrary base change.

\medskip\noindent
Let $Y = V_1 \cup \ldots V_n$ be an affine open covering.
Let $U_i = f^{-1}(V_i)$. If the statement is true for each of the
morphisms $U_i \to V_i$ over $R$, then it holds for $f$.
Namely, the scheme theoretic image of $U_i \to V_i$
is the intersection of $V_i$ with the scheme theoretic
image of $f : X \to Y$ by Morphisms, Lemma
\href{https://stacks.math.columbia.edu/tag/01R8}{lemma quasi compact scheme theoretic image (Tag 01R8)}.
Thus we may assume $Y$ is affine.

\medskip\noindent
Let $X = U_1 \cup \ldots U_n$ be an affine open covering.
Then the scheme theoretic image of $X \to Y$ is the same as
the scheme theoretic imge of $\coprod U_i \to Y$.
Thus we may assume $X$ is affine.

\medskip\noindent
Say $X = \Spec(A)$ and $Y = \Spec(B)$ and $f$ corresponds
to the $R$-algebra map $\varphi : A \to B$.
Then the scheme theoretic image of $f$ is $\Spec(A/\Ker(\varphi))$
and similarly after base change (by an affine morphism, but it is
enough to check for those).
Thus formation of the scheme theoretic image commutes
with base change if $\Ker(\varphi \otimes_R R') = \Ker(\varphi) \otimes_R R'$
for all ring maps $R \to R'$.

\medskip\noindent
After replacing $R$, $A$, $B$ by $R_g$, $A_g$, $B_g$ for a suitable
nonzero $g$ in $R$, we may assume $A$ and $B$ are flat over $R$.
By Lemma \href{https://stacks.math.columbia.edu/tag/0H40}{lemma cokernel flat (Tag 0H40)}
we may also assume $B/A$ is a flat $R$-module.
Then $0 \to \Ker(\varphi) \to A \to B \to B/A \to 0$
is an exact sequence of flat $R$-modules, which implies
the desired base change statement.
\end{proof}
\end{document}
